% \section{The ending of the YVi}\label{ENDOFTML}
The ending of a text, especially the \colophon s in its manuscripts, can give us various information about the text itself, its author, provenance, transmission, etc. This applies also to the YVi. Its closing stanzas have already attracted some scholarly attention.\footnote{See \citet[420]{rukmani:yoga1}.}
Unfortunately, the ending of the YVi has not been reproduced properly in the 1952 edition. %What they presented was almost fictional. They altered the order of verses and introduced quite different wordings in some places.

%Having access to accurate data is prerequisite for investigation. %It is regrettable that such an important part as the ending of the text was not properly reported in the 1952 edition, possibly leading to arguments without basis. 
Here is an attempt to give a better picture of the ending of the YVi. I will first present transcripts from the manuscripts \Tm\ and L\@.\footnote{I am not concerned with the remaining manuscripts, \Td, M or A here. We know them to be direct descendants of either \Tm\ or L\@. (See pp.~\pageref{sec:stemma}\,ff.) Unlike the other parts of this volume where editions of the YVi was presented, the purpose of this section is to discuss the transmission history of the YVi. It was useful to include readings in those modern manuscripts when the goal was to reach the most plausible reading of the YVi (in its intended form). Readings resulting from the intelligence and knowledge of the modern copyists, particularly that of M, are worth consulting and reporting.} Then I will present an edition, accompanied by a translation and annotations, of the part that is common to both manuscripts. Then I will proceed with editions of additional endings (the \colophon\ or colophon stanzas) that are unique to \Tm\ and L each. A note on the presentation of the ending in the 1952 and discussions on the layers in the endings will follow. 

\section{Transcripts}
Here are transcripts of \Tm\ and L toward the end. The commentary proper ends with \ldots\ \emph{vedāraṇyakavad ity om iti} \citep[369]{yvi}.\footnote{After this, both manuscripts have the decorative sign that appears when the commentary on a \sutra\ ends. Here the sign is transcribed as \textdb{❋}. Cf.~p.~\pageref{punctuation}.} 

Conventions generally follow those that have been adopted to report readings in the critical apparatus.\footnote{See pp.~\pageref{sec:symbols}\,ff.\ for more complete explanations.} The ones used here are: 
  \ccs\ and \cce\ to represent canceled elements; 
  \ins\ and \ine\ to represent elements added later;
  [ and ] to represent physically lost portions (if signs are partially lost but recoverable with confidence, the brackets enclose the reading and if signs are completely lost, they are represented by center dots (\mist{1}) whose number represents the estimated number of lost signs);
  \{\ and \}\ to represent modifications to the text that do not necessarily involve simple deletion and insertion.
  
\newcommand{\HTG}{\textdb{❋}}
\subsection{The end of \Tm}
\Tm\ ends in folio 125 recto (line numbers are in parentheses):
\begin{quote}
	(1) ntarasvarūpatatvanirddhāraṇena tantrāntarīyakatvasva\-rū\-pa\-ni\-rā\-ka\-ra\ccs n\cce \ins ṇ\ine ena ca kaivalya\lost{4}vedi[tavyā\mist{1}]\-ta\-[tro\mist{3}maṃgalārtthaḥ] prayuktaḥ parameśvara\-nā\-māṃ\-ki\-ta\-ma\-sta\-ka\-tayā \ins ta\ine\ccs s ta\cce ntrapra\-ca\ins y\ine ārttham śāntyarttham ve\-dā\-ra\-ṇya\-ka\-va\-d ity om iti \HTG\ %
	(2) oṃkāro yasya vaktā samacarata phalaiḥ ka\-rmma yasmād aśeṣa\-n niṣkarmakleśapāko gha\-[ṭ]a\-ya\-[ti] sakalaṃ yaḥ phalen\{a ā $\rightarrow$ ā\}%
	\footnote{Due to the fact that in the Malayalam script, the diacritical sign for the vowel \emph{ā} is written as an independent letter, what happens here is that the diacritical sign for \emph{ā} has replaced what appears to be an independent vowel \emph{ā}. So, the manuscript originally had \emph{phalena ā°} and was corrected to \emph{phalenā°}.}%
	kri\-yāṇā\ccs m īśvaro\cce \ins m īś\ine ā\-na\-m īśvaro ya sthiti\-bhavanidhanapra\-kri\-yā\-ṇāṃ vi\ccs p\cce \ins dh\ine ā\-tā ddhyā\-ya\-n naś śu\-kli\-mā\-naṃ vya\-pa\-nu\-da\-tu ta\-rāṃ kṛṣṇimānaṃ sa kṛṣṇaḥ a\-ne\-ka\-pha\-% 
	(3)ṇa\-ra\-tnau\-gha\-vi\-dyo\-ti\ccs X\cce%
	\footnote{I cannot decipher the sign marked to be cancelled. It appears to have a \emph{d} element as first part of the conjunct letter. The rest of the conjunct letter could have a \emph{c}, \emph{ṣ}, and/or \emph{m} element. It could have the vowel \emph{i}, but the circle-like part of the letter is too small for the usual way of marking a syllable with the \emph{i} vowel.} %
	dyu\-na\-bho\-di\-śe yo\-gī\-ndrā\-ya pha\-ṇī\-ndrā\-ya tat pa\-ta\-ñja\-la\-ye namaḥ va\-da\-nā\-hi\-ta\-pū\-rṇṇa\-ca\-ndra\-kaṃ gurum īśānam a\-bhū\-ti\-bhū\-ṣa\-ṇa\-m pra\-ṇa\-mā\-my a\-bhu\-ñjaṃ\-ga\-saṃ\-gra\-ha\ccs m\cce\ins ḥ\ine\ bha\-ga\-va\-tpā\-da\-m a\-pū\-rvva\-śaṃ\-ka\-ram \HTG\ go\-vi\-nda\-bha\-ga\-vatpū\-jya% 
	(4)pā\-da\-śi\-ṣya\-sya pa\-ra\-ma\-haṃ\-sa\-pari\-vrā\-ja\-kā\-cā\-ryya\-sya śaṃ\-ka\-ra\-bha\-ga\-va\-taḥ kṛtau pā\-ta\-ñja\-la\-yoga\-śā\-stra\-bhā\-ṣya\-vi\-va\-ra\-ṇe catu\-rtthaḥ pādaḥ \HTG\ yā\-trā\-kṛ\-tas tri\-ja\-ga\-tāṃ ya\-da\-ci\-ntya\-śa\-kti\-le\-śān da\-śā\-hur iha mī\-na-mu\-khā\-va\-tā\-rā\-[\ccs t\cce\ins n\ine] kle\-śo\-% 
	(5)\-dbha\-va\-tri\-vi\-dha\-tā\-va\-va\-ra\-mpa\-rā\-rttās tan nā\-ga\-nā\-tha\-śa\-ya\-naṃ śa\-ra\-ṇaṃ vra\-jā\-maḥ yo\-ge\-na ci\-tta\-sya pa\-de\-na vā\-cā\ins ṃ\ine ma\-la\-śa\-rī\-ra\-sya ca vai\-dya\-ke\-na yo pā\-ka\-ro\ccs kta\cce\ins t ta\ine m pra\-va\-ra\-mu\-nī\-nām \ccs v\cce%
	\footnote{This \emph{va} is clearly marked to be cancelled but I do not find the obvious replacement, \emph{pa}, in the manuscript. Also, the diacritical sign for \emph{ā}, too, should be cancelled, but I do not see a trace of such  cancellation.}%
	āta\-ñja\-li\-m prā\-ñja\-lir \ccs X\cce ā%
	\footnote{What appears to be written here is the sign \emph{ra}, followed by what appears to be the diacritical sign to mark the following consonant sign with the vowel \emph{e} (which is cancelled), and another such sign to mark the preceding consonant with the vowel \emph{ā}. What probably happened here was that the scribe started to write \emph{rane} or \emph{rano}, but realized the mistake and wrote \emph{rā na}.}%
	na\-to smi \HTG\ \HTG\ 
	(6) śrīmat pa\-ta\-ñja\-li\-ma\-ha\-rṣi\-vara\-pra\-ṇī\-ta\-\ccs ṃ\cce śrī\-yoga\-bhā\-ṣya\-hṛ\-da\-ya\-pra\-ka\-ṭī\-kṛd etat śrī\-śāṃ\-ka\-ra\ccs ṇa\cce\-ṃ vi\-va\-ra\-ṇaṃ sa\-ma\-va\{tti $\rightarrow$ rtti\}śa\-tro \ccs [g $\rightarrow$ rgg]ā\cce \ins gā\ine rggyo%
	\footnote{\label{GATORGG}There are two stages of corrections. First the original scribe wrote \emph{°śatro gārggyo} as it should read. Someone altered \emph{gā} to \emph{rggā} by adding strokes to the original sign for \emph{ga}. This made the word \emph{samavartiśatro} in the vocative case to the \emph{samavartiśatroḥ} in the genitive case. Then again the whole of the modified \emph{rggā} was cancelled and in the margin \emph{gā} was inserted. These two stages of corrections may have been done by the same person. All these alterations are not inked.} %
	vya\-lī\-li\-kha\-m aha\ccs tvi\cce \ins n tva\ine\ccs ñca\cce da\-vā\-pti\-kā\-maḥ \HTG\ \HTG\ \ccs va\cce\ins pa\ine ra\-kro\-ḍaḥ \HTG\ \HTG\ \HTG
\end{quote}

\subsection{The end of L}
And here is a transcript of L (folio 138 recto):
\begin{quote}
(4) \ldots\ tatroṃkāro maṃgalārtthaḥ prayuktaḥ parame\-śva\-ra\-nā\-mā\-ṃki\-ta\-mastakatayā tantrapracayārtthaṃ śāntyartthaṃ vā ve\-%
(5)\-dā\-ra\-ṇya\-kavad ity om iti \HTG\ oṃkāro yasya vaktā samacarata phalaiḥ karmma yasmād aśeṣan niṣkarmmakleśapāko gha\-ṭa\-ya\-ti sakalaṃ yaḥ phalena kriyāṇām īśānām īśvaro ya sthi\-%
(6)\-ti\-bha\-va\-ni\-dha\-na\-pra\-kri\-yā\-ṇāṃ vi\-dhā\-tā dhyāyan naḥ śuklimānam vyapanudatu tarām kṛṣṇimānam sa kṛṣṇaḥ pha\-la\-ra\-tnau\-gha\-vi\-dyo\-ti [dyu]na\-bho\-di\-śe yogendrāya phaṇīndrāya [ta\mist{1}ta]ñja\-%
(7)\-la\-ye namaḥ padanāhitapū\-rṇṇa\-ca\-ndra\-kaṃ gurum īśānam abhūtibhūṣaṇam praṇamāmy abhu\-ñja\-ṃga\-saṃ\-gra\-ham bha\-ga\-va\-tpā\-da\-pū\-rvva\-śaṃ\-ka\-ram \HTG\ go\-vi\-nda\-bha\-ga\-va\-tpū\-jya\-pā\-da\-śi\-ṣya\-[sya \mist{1}]\-ra\-ma\-haṃ%
(8)sa\-parivrājakācā\-ryya\-sya śaṃ\-ka\-ra\-bha\-ga\-va\-taḥ kṛtau pā\-ta\-ñja\-la\-yo\-ga\-śā\-stra\-bhā\-ṣya\-vi\-va\-ra\-ṇe ca\-tu\-rttha\-ḥ pādaḥ \HTG\ sa\-mā\-pta\-ñ cedaṃ vivaraṇam \HTG\ yātrākṛtas tri\-ja\-ga\-tāṃ ya\-da\-[ci]\-ntya\-%
(9)\-śa\-kti\-leśān daśāhur iha mīnamukhā[\mist{3}n] kleśo\-t\-bha\-va\-tri\-vi\-dha\-tā\-pa\-pa\-ra\-mparā\-rttās tan nā\-ga\-tnā\-tha\-śa\-ya\-naṃ śa\-ra\-ṇam pra\-jā\-maḥ 
yo\-ge\-na ci\-tta\-sya pa\-de\-na vācāmmalāśarīrasya tu %
(10) vai\-dya\-ke\-na yo pākarot tam pra\ccs vara\cce munīnām pata\-ñja\-li[m] prāñjalir ā\-na\-to smi|\footnote{As I write on p.~\pageref{punctuation}, hardly any punctuation is used in either \Tm\ or L\@. Here, however, L has what appears to be a \danda.} jitam pata\-ñja\-li\-muninā yena śreyo\-rtthi\-sā\-rttha\-sa\-syā\-nām vihito harṣas tā\-pa\-tri\-ta\-ya\-hṛtā %
(138v1) dha\-rmma\-me\-ghe\-na| yas ta\-tā\-na bha\-va\-mārgga\-laṃ\-ghi\-nām kle\-śa\-ka\-rmma\-ma\-ya\-gha\-rmma\-tu\-tta\-ye dha\-rmma\-megha\-mu\-kha\-yogatoya śrī\-ke\-śa\-va\-pra\-kā\-bha\-ṭṭā\-ra\-kā\-ṇāṃ śiṣyasya sa\-tyā\-na\-nda\-sya iya\-m pādañjalaṭī%
(2)kā \HTG\ \HTG\ \HTG
\end{quote}

\section{An edition of the common part}
Quite clearly, the ending is shared between the two manuscripts up to some point; the shared part must have been in the common ancestor,\footnote{Note that this does not necessarily mean that \emph{only} the shared part was in the common ancestor. There are possibilities that the common ancestor had more and that some of the continuation made it to either or neither of \Tm\ or L\@. It is, however, highly unlikely that the text shared in both \Tm\ and L originated independently.} which I consider to have been the exemplar of \Tm.\footnote{See section ``\titleref{extm}'' on pp.~\pageref{extm}\,ff.}  Here is an edition, translation of the common part along with some comments.

\subsection{The end of the commentary}
The YVi as a commentary ends as follows:
\begin{quote}
{\dn तत्रौंकारो मंगलार्त्थः प्रयुक्तः। परमेश्वरनामांकितमस्तकतया तन्त्रप्रचयार्त्थं शा\-न्त्य\-र्त्थं वा वेदारण्यकवदित्योमिति॥}

There the syllable \OM\ that has the auspicious meaning is used. [It is used,] since it is the foremost of things bearing the name of \Isvara, so that the system (\tantra) [expounded in this \sastra] grows, or for the sake of tranquility, analogous to the \Veda s and the \Aranyaka s [ending with the syllable \OM]. Therefore [the \sastra\ ends with] ``\OM.''/Thus \OM. End.
\end{quote}
This concludes the commentary. It appears that the author loaded several meanings in the last words of his text. The last part \emph{om iti} can be seen as the quote of the end of the root text, but it is at the same time the conclusion of his own text, too. In that sense, readers understand that he, too, placed the syllable \OM\ at the end in the sense he has just explained. This is the same technique the author uses at the beginning of the commentary. He starts the commentary with \emph{athetyādi pātañjalayogaśāstram. tasya vivaraṇam ārabhyate}.\footnote{See p.~\pageref{FIRSTSTC} for the text and p.~\pageref{ATHAATHATRL} for the translation.} He uses the very first word of the root text as the first word of the commentary. Also, the very last word, \emph{iti} serves several functions. One is to signal the end of a quotation from the root text. Another is as part of the quote of the very last word of the root text; we naturally expect the word \emph{iti} at the end of the root text. The third is to mark the end of the commentary itself. The part \emph{parameśvaranāmāṅkitamastakatayā} alludes to the statement in the commentary on \rYS{1.27}{1.27} where the syllable \OM\ is said to be the favorite name of \Isvara.\footnote{See p.~\pageref{FAVNAMEED} for the edition and p.~\pageref{FAVNAMETRL} for the translation, and n.~\ref{FAVNAME} on p.~\pageref{FAVNAME} for parallel passages in the \BAUBh\ and the \ChUBh.}

\subsection{First common concluding stanza}
\label{krsnaverse}
This conclusion is followed by a stanza in the \sragdhara\ meter:
\begin{verse}
%oṃkāro yasya vaktā samacarata phalaiḥ karmma yasmād aśeṣan 
%  niṣkarmmakleśapāko ghaṭayati sakalaṃ yaḥ phalena kriyāṇām| 
%īśānām īśvaro yaḥ sthitibhavanidhanaprakriyāṇāṃ vidhātā 
%  dhyāyannaḥ śuklimānam vyapanudatu tarāṃ kṛṣṇimānaṃ sa kṛṣṇaḥ|
\textdn{ओंकारो यस्य वक्ता समचरत फलैः कर्म्म यस्मादशेषन्निष्कर्म्म\kle शपाको घटयति सकलं यः फलेन क्रियाणाम्। \\
ईशानामीश्वरो य स्थितिभवनिधनप्रक्रियाणां विधाता 
  ध्यायन्नः शुक्लिमानं व्यपनुदतु तरां कृष्णिमानं स कृष्णः॥}\\
{\footnotesize pāda b: phalena kri° L] phalena ākri° \Tmac, phalenākri° \Tmpc\\
pāda c: īśānām īśvaro L] īśvaro \Tmac, īśānam īśvaro \Tmpc; vidhātā L\Tmpc] vipātā \Tmac

}
The syllable \OM\ signifies him; from him all the actions appeared along with [their] fruits; \\ He, [himself] having no actions, impurities or [their] ripening, unites everyone with the fruit of actions; \\
  He is the \inidx{lord} of lords; he performs the procedures of stability, rise, and end.\\ May that \Krsna, thinking of my good, utterly remove [my] wrong.
\end{verse}
This stanza applies a similar playful composition technique to that just seen in the end of the commentary. The author finished the commentary with the syllable \OM. Now a closing stanza immediately starts with the same word. This is deliberate and the mind behind this choice appears to be the same as that which chose to finish the commentary with the syllable \OM. This stanza is thus probably integral to the bulk of the commentary.

Much of the content in this stanza also agrees with what is put forward in the \Isvara\ section of the commentary and the opening stanzas.\footnote{Both are dealt with in this volume. For the opening stanzas, see pp.~\pageref{edition01}\,ff.\ for the text and pp.~\pageref{OSTZ1}\,ff.\ for translation. The \Isvara\ section is the main part of this volume.} The first segment of the stanza refers back to \rYS{1.27}{1.27}. The second and third segment (about \Isvara's being the legislator and enforcer of the law of \karman s) strongly resonate with the first \emph{pāda} of the first stanza in the opening of the whole text (see p.~\pageref{FVFPEDN} for the text and p.~\pageref{FVFPTRL} for its translation). It is interesting that \Isvara's being the creator of the laws of \karman, rather than his being the author of the Vedas that prescribe them, is more or less clearly stated here. In the opening stanza and in the \Isvara\ section of the commentary, this point was ambiguous. Also notable is the similarity in wording to that of the \BhGBh. See n.~\ref{NONFSFP} on p.~\pageref{NONFSFP} for the wording in the \BhGBh. In that connection, the explicit naming of \Isvara\ as \Krsna\ is also noteworthy. The author of this commentary repeatedly hints that he considers \Isvara\ to be \Visnu, especially the \Bhagavat\ of the \BhG\ (see pp.~\pageref{narayana}, \pageref{gita1155}, \pageref{omvisnu}, \pageref{NRYNN25}, \pageref{NNRYN26}, \pageref{VISNU27}, and \pageref{VSNSM1}), who indeed is \Krsna. The last \emph{pāda} further employs a word-play, between black/bad (\emph{kṛṣṇiman}) and \Krsna, alluding to terminologies in the text that has been commented; \rYS{4.07}{4.7}, for example, refers to actions as white or black (and their mixture).

In the last pāda the author asks \Isvara, who is revealed to be \Krsna, for his forgiveness. Mentioning \Isvara\ or \Krsna\ as someone who can remove bad actions attracts attention. I am not aware of any passage in the YVi that expresses the view that \Isvara\ can remove bad \karman. Hence this last \emph{pāda} sounds a little like being addressed to a human king, a judicial authority who, in reality, can grant pardons. %The author sounds as though soliciting \Krsna\ to be tolerant in judging actions. 

The whole stanza gives the impression that it is about a governing personage. This impression comes from the lack of the mention of \Isvara's being omniscient or omnipotent---the two aspects the author of the YVi usually mentions along with his being the lord of all. The last pāda, by introducing the view not put forward in the text, reinforces that impression. As a result, this stanza can be read as a prayer to \Krsna\ the deity as well as a plea to \Krsna\ a human ruler.\footnote{This point catches my attention because I propose the theory that \Sankara, the author of the \BSBh, lived under the rule of the Rāṣṭrakūṭa king, Kṛṣṇa I \citep{harimoto:2006}.} The author has already shown some literary techniques in this stanza. It would not be surprising if he meant this stanza to be read in these two meanings.

\subsection{Second common concluding stanza}
\label{patanjaliverseone}
Then \Tm\ and L records a śloka:
\begin{verse}
%anekaphaṇaratnaughavidyoti dyunabhodiśe 
%yogendrāya phaṇīndrāya tat patañjalaye namaḥ||

\textdn{अनेकफणरत्नौघविद्योतिद्युनभोदिशे। \\
योगीन्द्राय फणीन्द्राय तत्पतञ्जलये नमः॥}\\
{\footnotesize 
pāda a: anekaphaṇa° \Tm] phala° L\\
pāda c: yogīndrāya \Tm] yogendrāya L \\
pāda d: tat patañjalaye \Tm] [ta\mist{1}ta]ñjalaye L

}
To him who irradiate the heaven, sky, and directions with multitude of jewels residing in [each of] his many hoods \\
Thus salutation to [him,] the \inidx{lord} of yogins, the lord of serpents, \Patanjali.
\end{verse}
The above translation is tentative. As presented, I do not think the Sanskrit text of the stanza is intelligible. Yet I cannot offer a better solution to make the text meaningful. Either the transmission of this stanza was bad or it was not well composed in the first place.

In pāda a, I have adopted the reading \emph{anekaphaṇaratnaugha°} found in \Tm. L does not have \emph{aneka°} and hence the first line is short of three syllables there. The 1952 edition,\labelh{EDITALT0} based primarily on the manuscript M (an \inidx{apograph} of L) supplied \emph{pṛthivī°} in the beginning of pāda b, resulting in the reading \emph{phaṇa\-ratnau\-gha\-vidyoti [pṛthivī]dyunabhodiśe} for the first line.\footnote{L, and hence M, read \emph{phalaratnau°} rather than \emph{phaṇaratnau°}. That the 1952 edition still reads \emph{phaṇaratnau°} is an emendation.} Although I do not think this was a successful emendation, it was a clever one. When one sees \emph{dyaus} and \emph{nabhas}, it is natural to expect something that means ``earth.'' I do not adopt this reading because we do have a reading from \Tm\ that satisfies metrical requirements. In addition, the word \emph{diś} in the compound might have meant the same thing, the earth. 

The reading adopted here, \emph{anekaphaṇaratnaugha°}, is still not completely satisfactory because \emph{aneka} and \emph{ogha} seem semantically redundant. I cannot eliminate the possibility that the word \emph{aneka} in the beginning of the line was an emendation on the part of the scribe of \Tm\ (\Tm\ is an \inidx{apograph} of the common ancestor of \Tm\ and L). 

In pāda c, while L reads \emph{yogendrāya}, \Tm\ reads \emph{yogīndrāya}. I have adopted more common word \yogindra\ of \Tm. If we applied the principle of \emph{lectio difficilior}, perhaps \emph{yogendrāya} would be more likely to have been in the common ancestor. Given how syllables with the vowel \emph{e} is written in the Malayalam script, it is hard to imagine how \emph{gī} or \emph{ge} becomes the other by misreading. If, then, the change was intentional, it would be from \emph{ge} to \emph{gī} (\yogindra\ being more common Sanskrit word). Note, however, that being in the common ancestor does not necessarily mean that the stanza was originally composed with that reading. I have decided to adopt the word that would have been used by an author who knew standard Sanskrit; that is, I have followed the probable judgement of the scribe of \Tm. Nonetheless, as a result of this decision, the constituted reading is more or less that of \Tm. This is quite a contrast to the previous stanza where the scribe of \Tm\ was apparently struggling to record correct readings (see the variants for the previous stanza). If there is a trend that L was more successful in reproducing readings in the common ancestor, more acceptable readings in \Tm\ might in fact be a result of conscious alterations. Note that \emph{phaṇa\-ratnau\-gha\-vidyoti} that starts the stanza in L is in itself a good pāda of \sloka. And there is a sign that I cannot decipher in \Tm\ after \emph{°vidyoti}. Perhaps, despite metrically more correct or using more common words, the readings adopted from \Tm\ may still be far from what was originally intended.

I have further problems with this stanza. One of them is the word \emph{tat} in the second line. Again, a pronoun that points to a missing entity. It cannot be the first member of a compound (or can it?) affixed to a proper name. As an independent word, what we need is \emph{tasmai} that does not fit metrically. As the last resort, I have translated it as an adverb, meaning ``therefore.'' It still is preferable if the preceding line had a relative pronoun such as \emph{yat}, \emph{yatas} or \emph{yasmāt}. We cannot easily supply it as an emendation in the previous line.

It to me appears that this stanza is a remnant of badly transmitted text. It could have constituted of more than just 4 pādas. One possibility is that the original text was a sort of \emph{nirukti} of the name \Patanjali, trying to explain that it constitutes of elements of the verb \emph{pat-} (to fall) and \emph{jalin} (possessing water), perhaps in the sense that he rains down gems that shine the universe. The association between snakes and rain is a universal topos. 

\subsection{Third common concluding stanza}
\label{THETHIRDST}
Next we have a stanza in the \viyogini\ meter, dedicated to the first \Sankara:

%vadanāhitapūrṇṇacandrakaṃ gurum īśānam abhūtibhūṣaṇam|
%praṇamāmy abhujaṃgasaṃgraham bhagavatpādam apūrvvaśaṃkaram|| (Viyoginī)

\begin{verse}\label{SANKARASIVASTANZA}
\textdn{वदनाहितपूर्ण्णचन्द्रकं गुरुमीशानमभूतिभूषणम्।\\
प्रणमाम्यभुजंगसंग्रहम् भगवत्पादमपूर्व्वशंकरम्॥}\\
{\footnotesize 
pāda a: vadanāhita° \Tm] padanāhita° L\\
pāda c: abhujaṃga° L] abhuñjaṃga° \Tm; °saṃgraham L\Tmac] °saṃgrahaḥ \Tmpc \\
pāda d: °pādam apūrvva° \Tm] °pādapūrvva° L

}
The one to whose face the full moon is attached, the master, the ruler, the one not adorned by ashes\\
I salute him, without a collection of snakes, the Bhagavatpāda, the unprecedented \Sankara.
\end{verse}
This stanza is interesting in that it emphatically distinguishes \Sankara\ the teacher from \Siva, one of whose popular epithets is \Sankara. The first attribute mentioned here is the full-moon like face of \Sankara, a common description of a beautiful face. Its wording \emph{vadanāhitapūrṇacandra}, literally ``to whose face the full moon is attached'' is a contrast to \Siva, to whose head the crescent moon is attached. The expression \emph{abhūtibhūṣaṇam} appears to be a straightforward negation of \Siva's attribute of wearing ashes. However, there can be more positive interpretations such as ``not wearing any ornament,'' or ``to whom not owning wealth is an ornament.'' % the last one is Haru's suggestion.
This should be a positive attribute for a renouncer. Another contrast to \Siva, \emph{abhujaṃgasaṃgraha}, may be a play on the two major meanings of the word \emph{bhujaṃga}, the snake and the lover. Since the word \emph{saṃgraha} has a wide range of meanings, it is difficult to construe it with the second meaning. One possibility is that the compound means something like ``the one who does not partake in intercourse as a lover.'' This would be a description of again a renouncer. Nonetheless, it appears safe to assume that the technique of \emph{śleṣa} is used in these three attributes.

%The author of this stanza might have had a slightly sectarian attitude toward \Saivism. The attributes he says \Sankara\ does not have (wearing ashes and surrounded by snakes) represent the fearful aspect of the god. %; they are typically mentioned in stories when someone has negative things to say about \Siva. %[version 2] % this does not seem justifiable. % or does it?
%This is quite a contrast to hagiographical accounts or the modern popular perception that strongly associate \Sankara\ with \Saivism. 
%The author of this stanza at least wanted to assure that when he paid homage to “\Sankara,” it was directed only to the teacher and not (additionally) to the god.

%This emphatic discrimination between the god \Siva\ and the human teacher \Sankara, despite sharing the name, is quite a contrast to the late or modern tradition (or tendency?)\ to associate \Sankara\ with \Siva. This stanza even has a slight negative tone toward \Siva. Some of the attributes that \Sankara\ is said not to have (wearing ashes, surrounded by snakes) in this stanza represent horrifying aspects of the god. It does not appear that the author of this stanza was very sympathetic toward \Siva\ worshippers. The origin of this stanza may go back to the period when the cult of \Sankara\ was not yet strongly associated with \Saivism. [version 1]

A few more things may be noted. One is that \Sankara\ is referred to as a \inidx{lord} (\emph{īśāna}). Īśāna is another popular name of \Siva\ and this one is not negated. The application of this name to \Sankara\ indicates that his authority as a spiritual leader had already been established when this stanza was composed. In addition, the author of this stanza used the expression \emph{apūrva} (the first), which implies that there were succeeding \Sankara s. This stanza presupposes the tradition of \Sankaracarya s already in place. We do not know when the \Sankara\ \matha s where their heads are called \Sankaracarya s were institutionalized, but their earliest possible establishment is during \Sankara's lifetime. Hence this stanza must be from at least a few generations after \Sankara's life time. %Yet, this order had not associated the founder to \Siva. The verse reads as though the author tried to avoid confusing the founder of the order with \Siva, particularly negating the fearful aspect of the god. Given that earlier followers, possibly including contemporary disciples, pay homage to \Visnu\ \citep[153]{hacker:relations}, the provenance of this verse might not be so far from them.

\subsection{Colophon of the work}
\label{WORKCOL}
A sub-colophon for the fourth pāda and the \colophon\ for the entire commentary follow:
\begin{quote}
{\dn गोविन्दभगवत्पूज्यपादशिष्यस्य परमहंसपरिव्राजकाचार्य्यस्य शंकरभगवतः कृतौ पातञ्जलयोगशास्त्रभाष्यविवरणे चतुर्त्थः पादः ॥ समाप्तञ्चेदं विवरणम् ॥}
\end{quote}
The first part, the sub-colophon, is a fairly standard \colophon\ to a work ascribed to the famous \Sankara.\footnote{See \citet{hacker:avsbhkl}.} At least, whether true or not, there is no ambiguity to whom the work is ascribed: not a random person called \Sankara, but the author of the \BSBh.\footnote{Another context in which this \colophon\ is significant is the discussion on the title of the work. For that, see \index{Wezler, Albrecht}\citet[17]{wezler:yvi1}} What troubles me slightly is the fact that there has been no \emph{iti} preceding these colophons. The end of the text is not clearly marked. This may or may not be a sign that some of the preceding concluding stanzas were not by the author of the bulk of the text itself.

\subsection{Post-colophon concluding stanza 1}
The \colophon\ is followed by two more stanzas in both \Tm\ and L\@. This first one is as follows:
\begin{verse}
%yātrākṛtas trijagatāṃ yadacintyaśaktileśān daśāhur iha mīnamukhāvatārān|
%kleśodbhavatrividhatāpaparamparārttās tan nāganāthaśayanaṃ śaraṇaṃ vrajāmaḥ||
{\dn यात्राकृतस्त्रिजगतां यदचिन्त्यशक्तिलेशान्दशाहुरिह मीनमुखावतारान्।\\
\kle शोद्भवत्रिविधतापपरम्परार्त्तास्तन्नागनाथशयनं शरणं व्रजामः॥}

{\footnotesize 
pāda a: yadacintya° \Tm] yada[ci]ntya° L\\
pāda b: °mukhāvatārān \Tmpc] °mukhāvatārāt \Tmac, °mukhā[\mist{3}n] L\\
pāda c: °tāpaparamparārttās L] °tāvavaramparārttās \Tm\ \\
pāda d: nāganātha° \Tm] nāgatnātha° L; vrajāmaḥ \Tm] prajāmah L

}
They say, in this world, that the ten \avatara s, starting with the fish, that keep the three worlds going, are fractions of his inconceivable power. \\%``The ten \avatara s, starting with the fish, in this world are fractions of his inconceivable power,'' say those who make the rounds of the three worlds,\\
We, being continuously afflicted by three kinds of sufferings originated from the impurity, seek refuge in him, the one who sleeps on the serpent king.
\end{verse}
This stanza in the \vasantatilaka\ meter has few textual problems (mostly exchanges between \emph{p} and \emph{v}, typical of the Malayalam script). It shows, in more or less a straightforward manner, an allegiance to \Visnu. 
%The only slightly tricky part in interpretation is the first word \emph{yātrākṛtaḥ}. Combined with \emph{trijagatām}, I take it to mean those who repeat the cycle of re-incarnations, i.e., ordinary beings. This goes best with the first word in the third pāda (\emph{kleśo\-dbhava\-tri\-vidha\-tāpa\-para\-mparā\-rttāḥ}), in the same case. 
It should be noted that \emph{tan} in the text in fact is \emph{tam}, an inflected word. The consonant \emph{m} is homogenized with the following \emph{n}.\footnote{An unlikely, yet not impossible, interpretation is that the \emph{tan-} is read as part of a compound \emph{ta\-nnāga\-nātha\-śayanam} (``his [\Visnu's] serpent king bed,'' i.e., \Sesa). If that (too) is meant, then the speaker of the stanza seeks refuge in \Sesa\ while praising \Visnu. Note that \Sesa\ and \Patanjali\ are sometimes identified.}
\subsection{Post-colophon concluding stanza 2}
\begin{verse}
{\dn योगेन चित्तस्य पदेन वाचाम्मलं शरीरस्य च वैद्यकेन।\\
यो ऽपाकरोत्तम्प्रवरम्मुनीनाम्पतञ्जलिम्प्राञ्जलिरानतो ऽस्मि॥} % upajāti

{\footnotesize 
pāda a: vācām \Tmpc (vācāṃ) L] vācā \Tmac\\
pāda b: malaṃ \emph{em.}]  mala° \Tm, malā° L; ca \Tm] tu L\\
pāda c: pākarot tam \Tmpc L] pākaroktam \Tmac; pravara \Tm\Lac] pra\ccs vara\cce\ L \\
pāda d: patañjalim L] \ccs v\cce ātañjalim \Tm, patañjali[m] L; °jalir ānato \Tmpc L] jalir \ccs X\cce ānato \Tm

}
The one who removed the impurity from the mind by yoga, from the speech by word, and from the body by medicine, \\
the best of the sages, \Patanjali, I bow to him, with my hands forming \emph{prāñjali}.
\end{verse}
This is, as \citet[(22)]{endoh:notes} points out, a stanza widely known among \inidx{grammarian}s.\footnote{See also \citet[xiv]{woods:hos} and \citet[503]{mbh1}. Cf.\ \citet[141--4]{meulenbeld:himia}} As such, it is doubtful that  this stanza was part of the original composition of the YVi.

This is the end of the shared text between \Tm\ and L\@. Following this stanza, \Tm\ has another \colophon\ stanza.

\section{An edition of the colophon stanza of T\kern -0.15em\raisebox{-.4ex}{\scriptsize m}}\label{TMCOL0405}
The \colophon\ stanza of \Tm\ involves many corrections. The text that was meant can be reconstructed as follows:
\begin{verse}
{\dn श्रीमत्पतञ्जलिमहर्षिवरप्रणीतश्रीयोगभाष्यहृदयप्रकटीकृदेतत्।\\
	श्रीशांकरं विवरणं समवर्त्तिशत्रो गार्ग्यो व्यलीलिखमहन्त्वदवाप्तिकामः॥} % vasantatilakā

{\footnotesize 
pāda a: °praṇīta° \Tmpc] °praṇītaṃ \Tmac\\
pāda c: śrīśāṃkaraṃ \Tmpc] śrīśāṃkaraṇaṃ \Tmac; 
        samavartti° \Tmpc] samavatti° \Tmac\\
pāda c-d: °śatro gārggyo \Tmac]  śatro \ccs [g $\rightarrow$ rgg]ā\cce \ins gā\ine rggyo \Tm\ (see note \ref{GATORGG})\\
pāda d: ahan tvad° \Tmpc] aha\ccs tvi\cce \ins n tva\ine\ccs ñca\cce d° \Tm

}

This [text] reveals the essence of the Yogabhāṣya composed by the venerable \Patanjali, the best of the great sages \\
I, \Gargya, have made it written, the commentary by venerable \Sankara, o slayer of \Samavartin, hoping to reach you.
\end{verse}
This stanza is followed by the very last word in the manuscript:
\begin{quote}
{\dn परक्रोडः}
\end{quote}
of which the first letter is originally written \emph{va} and corrected by a later hand (not inked) to \emph{pa}.

This \colophon\ informs us of several things. First, we learn the name, \Gargya, of the person who was one way or another responsible for the production of \Tm. He shows allegiance to \Samavartisatru, possibly \Visnu\ or \Siva: the slayer of \Samavartin\ (\Yama).\footnote{Dictionaries list names such as Yamāri, Yamaghna, Yamaripu as names of \Visnu. I am not aware of any work in which such names occur as such. On the other hand, a name similar in meaning, \Yamantaka\ is well-attested as a name of \Siva.} This much is certain.\footnote{Another potentially interesting point exists. The person who wrote this stanza considers \Patanjali\ to be the author of the \Yogabhasya.} 

However, there are many uncertainties. I am not certain if the name \Gargya\ belongs to the individual or it is his \gotra{} name. (For the possibility of it being a \gotra{} name, see below.) Nor am I certain in what capacity he was involved in the production of \Tm. The verb \emph{vi-likh-}, used in causative, usually means that the subject ordered the manuscript to be prepared. Most of the time, there is a mention of another name as the scribe.\footnote{These statements come from the observation of colophons of Nepalese manuscripts to which I have a good access to.}  But we do not have another name. We cannot take the name \Samavartisatru\ as that of the scribe. The word \emph{tvat-} “you” and the vocative case of \emph{samavartiśatro} compliment each other. Another curiosity is that the verb is in the first person. As far as the Nepalese manuscripts---to whose \colophon\ data I have an easier access---are concerned, when the verb \emph{vi-likh-} is used in the causative aorist, it is in the third person. The third person makes sense since the subject is not the person who was doing the actual writing. The stanza reads as though someone gave the scribe the text to be copied and the scribe reproduced it. I do not have enough experience to know if this was a common practice in Kerala.
%
%I cannot justifiably propose a better reading of the stanza to eliminate these oddities. %Yet, apart from some ambiguities, overall intention of the stanza presented here is more or less clear. 

It may further be noted that the constituted text of the stanza essentially follows the text after corrections in order to be more or less intelligible. The fact that there are so many corrections in the manuscript and the nature of the errors are also curious. 

Two of the errors originally made in the manuscript (\emph{praṇītaṃ} for \emph{praṇīta°} and \emph{śrīsāṃkaraṇaṃ} for \emph{śrīśāṃkaraṃ}) do not tell much of a story because they can occur in various circumstances;\footnote{They are surely mistakes. If the word \emph{śrīmatpatañjalimaharṣivarapraṇīta} ended there, then the \emph{etat} \ldots\ \emph{vivaraṇam} would be a work by \Patanjali. That would contradict with the adjective \emph{śrīśāṃkaram}. The word \emph{śrīśāṃkaraṇam} is unintelligible and would make the stanza unmetrical.} they are caused by mental expectations. Such writing mistakes could happen even when the person writing the manuscript was the composer himself.\footnote{As I see more colophons, I notice many mistakes are made in them even though the scribes were presumably writing their own words. I used to find it odd and suspected that colophons were being copied. However, it appears more likely that scribes are more used to copying than composing. Copying is a mechanical activity that one can train themselves to excel. However, composing and writing down at the same time is quite a different activity. I experience myself making more mistakes when writing something down especially by hand, at the same time thinking what to write. I find it not so rare to write my own name wrongly. I do not assume a \colophon\ being copied when mistakes, including the ones involving the scribe's own name, are found any more.} 
%
However, the mistake that was corrected to \emph{ahan tvad°} is rather inexplicable for any reason other than copying. The two signs of the original reading \emph{tvi} and \emph{ñca}, later replaced by \emph{ntva}, are completely unintelligible and hyper-metrical when read together. The only feature I find in these three signs is their graphical similarities. Unintelligible yet graphically similar signs appearing in a manuscript typically occurs when the scribe did not understand the nature of a correction. In this case, what happened was probably as follows: there was a manuscript where the part was first written \emph{vyalīlikham ahañ ca} and subsequently corrected to \emph{ahan tva°} by cancelling \emph{ñca} and inserting \emph{tva} in the margin; then a scribe (of \Tm) who was copying from this written text did not realize the cancellation and wrote down both, additionally misreading \emph{ntva} as \emph{tvi}. After \Tm\ was inked, someone further corrected it to the intended reading. The nature of errors being graphical, we can exclude the possibility that this error was produced by mishearing or mental associations. We can hence exclude the possibility that this \colophon\ stanza was being dictated (cf.\ the causative \emph{vyalīlikham} discussed above). 

Furthermore, I doubt that this stanza was present in the exemplar of \Tm. This is the very end of the manuscript where, in most cases, information relevant to the manuscript is recorded. There still is enough space on the side for about three lines of text and the side is the recto. It is not that the last folio was lost or that the scribe did not have enough space to write further. This stanza was intended to be the \colophon\ of the manuscript. The most likely scenario appears to be that it was prepared by the person who commissioned the manuscript (Gārgya) in written form and the scribe copied it.

I would now like to discuss the final word of the whole manuscript, \emph{pa\-ra\-kro\-ḍaḥ}.\footnote{I owe Yasutaka Muroya for calling my attention to the fact that \Parakroda\ is a name.} \Parakroda\ is the name of the family known to have produced \Jyesthadeva{} (active between 1500 and 1610)---the author of an astronomical work, \Yuktibhasa.\footnote{See \citet[187--188, and 207]{kvsarma:1991}. In addition to the post-colophon statements in a manuscript of a Malayalam commentary on the \Suryasiddhanta, Sarma cites Parayil Raman Namputiri, \emph{Nampūtirimār} (in Malayalam), Trichur 1918, p.~55, as a source of the information that \Jyesthadeva\ hailed from the family \Parakroda\ (Sanskritised name of Malayalam \emph{Paraṅṅoṭṭu}).} But \Parakroda\ may have been the name of the family deity as well.\footnote{This comes from the opening of a manuscript of the \Nyayabhasya\ to which Yasutaka Muroya drew my attention. The scribe of the manuscript shows allegiance to \emph{parakroḍadeva}. I do not have the details of the manuscript. Muroya kindly sent me a digital reproduction of the first side of the manuscript, and it starts with \emph{hari śrīgaṇapataye namaḥ śrīparakroḍadevāya namo namaḥ}. The first word here is \emph{hariḥ} whose \visarga\ has dropped before the combination of a sibilant and a semi-vowel.} I do not have a clear idea what function the word \emph{parakroḍaḥ} in the nominative singular at the end of \Tm\ serves. It would seem that it is there as an auspicious word.\footnote{This use would be analogous to the Nyāyabhāṣya manuscript mentioned in the previous note.} Yet again, the original scribe wrote \emph{varakroḍaḥ} instead of \emph{para°}. This is a typical error stemming from the graphical similarity between \emph{pa} and \emph{va} in the Malayalam script. As I have speculated above, the scribe might have been given what to write at the end of the manuscript and was copying it. He was probably not a \Parakroda\ himself; he was not familiar with the word.

It is in this connection that I cannot take the name \Gargya\ simply as a personal name. It can instead be the \gotra{} name of the person who commissioned \Tm. \Jyesthadeva\ mentioned above had some close relationship with another astronomer from Kerala, \NilakanthaSomayaji\ (born 1444 and composed the Tantrasaṃgraha, in 1500).\footnote{See \citet[187–8]{kvsarma:1991}. For the date of \NilakanthaSomayaji, see \citet[S 3]{sarmaandnarasimhan:tantrasangraha}.} They are said to have been fellow disciples of the same teacher and \Jyesthadeva\ based his \Yuktibhasa\ on the Tantrasaṃgraha. And this \NilakanthaSomayaji's \gotra{} was \Gargya.\footnote{See \citet[S 1]{sarmaandnarasimhan:tantrasangraha}.} Given the historically recorded relationship of a Gārgya and someone from the Parakroḍa family, it is hard not to consider the person referred to as \emph{gārgyaḥ} in \Tm's \colophon\ stanza as a \Gargya\ as well. On the other hand, it does not appear that \Jyesthadeva\ and \NilakanthaSomayaji\ came from the same family, and hence there is no need to consider that \Jyesthadeva's \gotra{}, too, was \Gargya, i.e., \Parakroda\ family's \gotra{} was \Gargya. I cannot deny the possibility that the nominative singular of the word \emph{parakroḍaḥ} might correspond to that of the word \emph{gārgyaḥ} in the \colophon\ stanza; but connecting these two words is not necessary.

There is too little information in the \colophon\ stanza to determine the provenance of \Tm\ precisely. The most we can say is that it may have had something to do with the circle of Kerala astronomers of the late 15th to early 16th century.


\section{An edition of colophon stanzas of L}\label{LCOL0405}
L has other dedicatory stanzas and a brief colophon. The text and a tentative translation of stanzas may be presented as follows:
%jitam pata\-ñja\-li\-muninā yena śreyortthisārtthasasyānām vihito harṣas tāpatritayahṛtā%
%(138v1)dharmmameghena| yas tatāna bhava\-mārgga\-laṃ\-ghi\-nām kleśakarmmamayagharmmatuttaye dharmma\-megha\-mu\-kha\-yogatoya śrīkeśavaprakābhaṭṭārakāṇāṃ śiṣyasya satyānandasya iya\-m pādañjalaṭī%
%(2)kā \HTG\ \HTG\ \HTG
 
%jitam patañjalimuninā yena śreyortthisārtthasasyānām| \\ % 12 18
%vihito harṣas tāpatritayahṛtā dharmmameghena|| % 12 15
\begin{verse}
{\dn जितम् पतञ्जलिमुनिना येन श्रेयोर्त्थिसार्त्थसस्यानाम्। \\
विहितो हर्षस्तापत्रितयहृता धर्म्ममेघेन॥}

Victorious is the sage \Patanjali, who caused joy by means of the cloud of Dharma (the \Dharmamegha [\samadhi]), the remover of heat (three torments), to the crops which are the multitude of those who seek better fortune.

%yas tatāna bhavamārggalaṃghinām % rathoddhatā
%kleśakarmmamayagharmmanuttaye|
%dharmmameghamukhayogatoyadaḥ
%śrī [patañjalimuniṃ taṃ astu śrīḥ]||


{\dn यस्ततान भवमार्ग्गलंघिनां
\kle शकर्म्ममयघर्म्मनुत्तये। \\
धर्म्ममेघमुखयोगतोय\ldots}

[[I salute \Patanjali\ldots]] who expanded [[\ldots\ that produces]] water that is yoga [pouring] from the \Dharmamegha\ (Dharma-cloud) in order to remove the heat, consisting of impurities and deeds, from those who are traversing the path of reincarnation.
\end{verse}
As can be seen, I did not make any emendation from the reading found in the manuscript. It does appear that the scribe wrote down what was meant to be written down. However, there are problems that raise questions about the circumstances under which these two stanzas were recorded in this manuscript.

The first and most glaring defect is that the second stanza is not complete. It lacks 12 syllables at the end. The stanza is interrupted in the middle of a word and the scribe simply continues to write the name of the owner (not certain; see below) after \emph{°yogatoya}. The scribe did not leave any indication that something was missing.

Another outward problem is a metrical defect in the first stanza. It starts with a \emph{ja}-\gana\ ({\applesymbols ⏑\,–\,⏑}). This is a metrical defect in the \arya\ meter.%meant to be in the \arya\ meter, has a metrical defect. It starts with syncopation, in other words, the first \gana\ is a \emph{ja}-\gana\ ({\applesymbols ⏑\,–\,⏑}). An \arya\ stanza should not have that rhythm in the odd numbered \gana s (the first to third syllables, 7th to 9th syllables, and so on).
\footnote{Harunaga Isaacson pointed this out to me.}
% The second stanza lacks of 12 syllables at the end. We could give the author(s) of the stanzas the benefit of the doubt: transmission errors (which I have just said is unlikely). Still, the stanzas are not well composed. 
%
This\labelh{EDITALT1} was probably the reason why the editors of the 1952 edition silently emended the first \pada\ of this stanza to \emph{sa jayati patañjalimuniḥ} \citep[370]{yvi}, which means the same thing, is metrically sound, and has the added benefit of having the corresponding pronoun to the relative pronoun \emph{yena}, the first word of the second \pada. 

The first stanza also appears clumsy for having so many words in the instrumental singular, masculine (\emph{patañjalinā}, \emph{yena}, \emph{dharmameghena}, and \emph{tāpatritayahṛtā}). The first three each occupy different syntactic functions and the last modifies the fourth. In particular, I find the composition ``\emph{yena \ldots\ harṣo vihito \ldots\ dharmamethena}'' (the words rearranged for clarity) to mean ``who caused the joy be means of \emph{dharmamegha}'' not clean. There the words \emph{yena} and \emph{dharmameghena} refer to two different things in the same clause: the agent of the action and an instrument, respectively.
%The first \arya\ stanza has three things in the instrumental singular, masculine, that serve different functions, viz., \emph{patañjalinā}, \emph{yena} and \emph{dharmameghena} (\emph{tāpatritayahṛtā} modifies this). The first, \emph{patañjalinā}, functions as the correlative of \emph{yena}. An explicit use of the correlative (\emph{tena}) is omitted, which in itself is fine, but \emph{patañjlinā} does not have to be in the instrumental case, somewhat confusingly. It can be in any case. This use of impersonal passive construction bothered the editors of the 1952 edition so much that they silently rewrote the first pāda into \emph{sa jayati patañjalimuniḥ} \citep[370]{yvi}. Furthermore, \emph{yena} and \emph{dharmameghena} refer to two different things in the same clause: the agent of the action and an instrument, respectively. Again, it is not hard to understand what is meant, but not an elegant composition. 

The compound \emph{śreyorthisārthasasyānām} also looks somewhat clumsy. What the author of the stanza wanted to convey was probably that \Patanjali\ promised joy to the \emph{śreyo\-rthin}s in the form of crops or the fruits achieved by the heat-reducing, i.e., rain producing, cloud, viz., the \Dharmamegha\ \samadhi. The author of the stanza is apparently playing on the word \emph{megha} (cloud) part of the \Dharmamegha\ and the meanings of the word \emph{tāpa} (heat and suffering). The topos behind this is the autumnal cloud that produces rain, relieving people of the summer heat, and contributing to  agriculture. However, with the compound \emph{śreyo\-rthi\-sārtha\-sasyānām} ending with the genitive plural, followed by \emph{harṣa}, the natural interpretation is that the people are compared to crops and they experience joy. This probably is not what was meant. I would think people rejoice due to the removal of heat and the crops, both given by the autumnal cloud. %This interpretation still results in a clumsy sense because it is unclear whether the joy comes from the success (crops) or the fact that people are relieved from the heat/sufferings. 
I would attribute the clumsiness to the inability of the author of this stanza to formulate his idea precisely.

The second stanza in the \rathoddhata\ meter is, as has been mentioned, short of one syllable and one pāda.\footnote{\citet[370]{yvi} supplies \emph{°dam [patañjalim ṛṣiṃ praṇato ‘smi]}. Although\labelh{EDITALT2} the editors indicate that only the last pāda was their reconstruction by brackets, they in fact added \emph{°dam} to have \emph{°toyadam}.} The parts enclosed in the double brackets in the translation are what I conjecture to have been in the missing part. Interestingly, the topos of the stanza significantly overlaps with that of the previous one. This stanza is about \Patanjali\ who rescues those who are suffering from the karmic existence by means of yoga; the suffering is compared to the (summer) heat; and the \Dharmamegha\ \samadhi\ is compared to the (autumnal) cloud that produces rain. Seeing the word \emph{tatāna} ``expanded,'' I would speculate that there was also an imagery of the serpent king \Patanjali\ who expands the hood to cause rain. However, as the editors of the 1952 edition conjectured, the most fitting syllable to supply at the end of the c-pāda indeed is \emph{da(m)}. There are not many words that have only one syllable after \emph{toya}. But then the whole compound (\emph{dharmameghamukhayogatoyada}) becomes rather unintelligible. The relative clause would then mean something like ``the one who spread the cloud that is yoga where the Dharmamegha is the foremost.'' It is not impossible but not good. In addition, there is no sign that the omission in the manuscript was a copying mistake. The best explanation for the lack of 12 syllables appears to be that the scribe (author of the stanzas himself?) just gave up finishing the stanza. That the last two stanzas in L convey more or less the same content contributes to this suspicion. It is as if the author of these stanzas were experimenting with different expressions to convey the same sense and gave up on finishing the second attempt.

The puzzling conclusion of L continues with the final statement. Here is the text without any corrections:\labelh{CFSGCLP1}

\begin{quote}
\emph{śrīkeśavaprakā\emph{[sic]}bhaṭṭārakāṇāṃ śiṣyasya satyānandasya i\emph{[sic]}yam pāda\emph{[sic]}\-ñja\-la\-ṭīkā} 

(This [manuscript of the] Pātañjaliṭīkā belongs to \Satyananda, a disciple of Keśavaprakāśabhaṭṭāraka).
\end{quote}
I note three problems (one is minor). First of all, the teacher's name should be \emph{śrīkeśavaprakāśa\-bhaṭṭāraka}. The manuscript omits \emph{śa} in \emph{°prakāśa°}. Also, it misspells \emph{pādañjala°} for \emph{pātañjala°}. This is most likely an influence from the scribe's native language (presumably Malayalam) where a consonant is voiced or unvoiced according to the position. Another minor issue is the lack of the \sandhi\ between \emph{satyānandasya} and \emph{iyam}; if the \sandhi\ is applied, the text would read \emph{°nandasyeyam}. I am not certain if \Satyananda\ was the scribe or the owner of L\@. (It is also possible that he was the scribe and the owner.) At least, we can gather that he was closely involved with L and his teacher was Keśavaprakāśa. I consider this \colophon\ to be genuine, as opposed to being copied from elsewhere. He was thinking of the text word by word (lack of the \sandhi), he wrote down what he would pronounce rather than what he saw (\emph{pādañjala} rather than \emph{pātañjala}). He might have been losing focus and was in a hurry with the prospect of finishing the manuscript (dropping a syllable from the teacher's name).

\section{Notes on the 1952 edition}
As has been noted, in the closing part of the YVi the editors of the 1952 edition introduced some changes to the text, mostly silently.\footnote{See pp.~\pageref{EDITALT0}, \pageref{EDITALT1} and \pageref{EDITALT2}.} 
%Without access to manuscripts, the reader would not know what they did. The editors 
They made another, more significant, change to the text without telling the reader: they moved around the stanzas and the colophon.

According to the two principle manuscripts, the end of the YVi consists of the following elements. First, both have:
\begin{enumerate}
\item a \sragdhara\ stanza, a prayer to \Krsna
\item a (corrupt) śloka(?) paying homage to \Patanjali
\item a \viyogini\ stanza paying homage to \Sankara
\item \colophon\ to the work
\item a \vasantatilaka\ stanza expressing allegiance to \Visnu, and 
\item an \upajati\ stanza paying homage to \Patanjali
\end{enumerate}
then \Tm\ has the \colophon\ stanza. L further has:
\begin{enumerate}
\setcounter{enumi}{6}
\item a defective \arya\ stanza saluting \Patanjali
\item an incomplete \rathoddhata\ stanza (presumably) saluting \Patanjali
\item the ownership statement
\end{enumerate}
The 1952 edition ends its text with most of the above but in a different order: 1, 5, 2, 6, 7, 8, 3, and 4 (the colophon).\footnote{It does not have 9; this is understandable because it pertains only to the manuscript L\@.} The edition derives from the manuscript M, which is a copy of L,\footnote{See pp.~\pageref{l2m2a}\,ff.} and M has the same ending as L\@. So, it was the editors of the 1952 edition who rearranged the text. The result was that the YVi ends with: two stanzas (1 and 5) dedicated to \Visnu, four stanzas (2, 6, 7 and 8) to \Patanjali, one stanza (3) to \Sankara, and the colophon (4). We can easily infer the motivation behind the change. Apparently the editors were not happy with the original arrangement because the objects of dedication went back and forth; usually, stanzas dedicated to one object appear in succession. The editors tried to \emph{fix} this perceived peculiarity. They did not, however, take into account if there is any reasonable explanation how the arrangement they imagined could have been mangled so badly. I doubt there is any. They in fact assumed that all the concluding stanzas were part of the text and found the arrangement strange. I, on the other hand, find the assumption not warranted. 
% WThis would be true if a same person was paying homage to various objects.  I certainly cannot explain the process and I do not think the rearrangement is necessary. Rather, it is not called for. %account for was how one could explain the peculiarity occurring under the same assumption. They should have reflected on their assumption.%but there is no need to assume that. Moreover, I cannot conceive how the ending in the 1952 edition would have changed to the one found in the manuscripts. 

\section{Layers in the ending of the YVi manuscripts}
I rather view the seemingly peculiar arrangement of dedicatory stanzas in the end of the YVi as a sign of the following: the ending grew as manuscripts were copied because different scribes added their own dedicatory/colophon material without (completely) removing their predecessors'. 
%I rather find the arrangement of stanzas in manuscripts to be a sign that the ending of the text grew in size by adding more elements to the end. 
There are several other signs that the ending of the YVi manuscripts consists of several layers.

%It is a norm that Sanskrit manuscripts convey more than the text for which they were produced. The additional material may be of two kinds: the one that pertains to the text (title, author, etc.) and the other that pertains to the manuscript (various information about the scribe, copying circumstances). may be divided roughly into the one about the text \ and the other about the manuscript . It is another norm that scribes want to start and end their manuscripts with auspicious symbols, phrases (both the beginning and the end) or even verses (in the end). Sometimes the borders between the text proper, colophons, remarks about the text itself and those about the manuscript and their provenances (e.g., did the author write the colophons between chapters or did someone else write them?), are not clear. It can happen that some portion of the additional material for a particular manuscript is coped in a subsequent manuscript. If such a process repeats, succeeding manuscripts of the same text may have a bigger and bigger ending.\footnote{This does not seem to apply to the beginning of manuscripts. The format for the beginning of manuscripts has little variations.} The ending of the YVi manuscript appears to be such a case. Luckily for us, the process left some marks.
%Primarily, what we find in what order at the end of the YVi manuscripts It is evident that the ending of the YVi manuscripts had layers that indicate the transmission process (adding more to the end) of the text.

%The outermost layer in the ending of the YVi manuscripts is relatively clearly visible because 

One immediately sees that the texts of \Tm\ and L deviate at one point. That is after the item numbered 6, using the list in the previous section. This tells us several things.
%This is an indication that anything that follows 6 in \Tm\ or L were introduced somewhere, in other words, they are not essential part of the text. 
The first is that items 1 to 6 were in the latest common ancestor of the two manuscripts, $\alpha$ (\Tm's exemplar).\footnote{See pp.~\pageref{extmandl}\,ff.}

As far as \Tm\ is concerned, it is most likely that the colophon stanza of \Tm\ is unique to it. Similarly, the ownership statement of L (item 9) is probably unique to L\@. The two faulty stanzas (7 and 8) of L may well belong to the same layer, judging from the curious quality (full of errors) shared by the stanzas and the ownership statement. Even if they were in $\alpha$,\footnote{There are various possibilities with regard to the origin of the two stanzas. It is not inconceivable that stanzas 7 and 8, in their entirety or partially, were in $\alpha$ and the scribe of \Tm\ dropped the part thinking that it pertained only to his exemplar and not something he had to reproduce. The faulty nature of those stanzas, if they were in the exemplar, might have contributed to the decision. Also, we cannot preclude that items 7 and 8 were added during the transmission between $\alpha$ and L since there is a possibility that there were intervening manuscripts between $\alpha$ and L (see p.~\pageref{LandALPHA}). Other possibilities include: that the colophon stanza found in \Tm\ was in $\alpha$ but it was replaced in L while it survived in \Tm; and that $\alpha$ had more text that was dropped both in \Tm\ and L\@.} 
%We should still be careful. The exact provenance of the last two stanzas could be a little complex. 
%
%We cannot, however, state that $\alpha$ had only up to item 6.   Nonetheless 
the scribe of \Tm\ felt them auxiliary enough to drop them. We can, therefore, draw a border after stanza 6.

%To summarize, items 7 and 8 in the above list may be introduced at the end of L by the scribe. Yet there are some possibilities that they were, at once or gradually, introduced to the end of the YVi text transmission in any of the ancestors of L, including $\alpha$. Even if any part of those two stanzas was in $\alpha$, . The colophon stanza of \Tm\ and the ownership statement of L (item 9) were introduced by the respective scribe.

Another border is after the colophon (4), revealing a layer consisting of stanzas 5 and 6. Unless there is a compelling reason to suggest otherwise, I consider anything after the colophon to have been added by a scribe. The scribe does not have to be the one who produced the manuscript at hand. Even if later scribes reproduced what an earlier scribe had added, it is still by a scribe. What we have here appears to be such a case. I find no reason to consider the two stanzas not to be scribal. Stanzas 5 and 6 express allegiance to \Visnu\ and pay homage to \Patanjali\ respectively. This is consistent with what a scribe would do. He was expressing allegiance to his tutelary deity and paying homage to the author of the \PYS\ on which the YVi is a commentary. %This is why the homages to them appear to be interrupted in the series of stanzas at the end of the YVi manuscripts. There is no reason to move these stanzas around.

%They were probably additions during the transmission by a scribe because they follow the colophon. In general, anything that comes after the phrase \emph{samāptam \ldots} in manuscripts is very fluid. I usually consider the text there only relevant to the manuscript unless there is evidence suggesting otherwise, in this case the fact that some of the following text is shared in two manuscripts. A scribe may add information about himself and/or the manuscript, or he might write auspicious stanzas. He may replace what was in his exemplar, or he may reproduce some part of it. This probably is such a case. Auspicious stanzas were reproduced. %In this case, however, at some point in the transmission, even as late as in $\alpha$ (\Tm's exemplar), the scribe copied the two stanzas that were meant to be unique in his exemplar.
%\footnote{There might have been more , but we have no means to tell.}
%That the homage to \Patanjali\ is interrupted by several stanzas, that the homage to \Visnu\ again (although the first one was to \Krsna) (after stanzas 1 and 2), is an indication that they were additions by a scribe. As the editors of the 1952 felt, an integral text is unlikely to go back and forth between \Visnu\ and \Patanjali\ as the object of homage. Furthermore, stanza 6 is not original but a widely circulated one. Such a stanza is fitting for a scribe to add at the end of his manuscript. This might indeed have been the last text in a manuscript, baring short statements about the owner, such as the one found in L\@. %A scribe may use such a cookie-cutter stanza but probably not an author, unless he was the first to compose the stanza that later became widely used. %Such a stanza can be easily added at the end of a manuscript by a scribe who wants to make his manuscript a bit more auspicious.

%Whether there were more layers before the colophon in the ending in the YVi transmission is unclear. 
Now we are left with three stanzas (items 1, 2, and 3) and the colophon. I do not see any clear borders in this part. Here I start examination from the first stanza.

The first stanza, a prayer to \Krsna, appears integral to the body of the text as has been discussed.\footnote{See pp.~\pageref{krsnaverse}--\pageref{patanjaliverseone}.} 

I do not see any strong positive or negative indications whether the second concluding stanza (numbered 2 above)\footnote{See p.~\pageref{patanjaliverseone}.} is part of the original composition. One thing that attracts my attention is that it pays homage to \Patanjali. The reason why \Patanjali\ should be paid homage to in the context of a commentary on the \PYS\ is that he is considered the author of it. However, while the author of the YVi mentions the name \Patanjali, he shows no awareness that \Patanjali\ is considered the author of the work he is commenting on. \Patanjali\ is referred to only as the author of the \VMBh, a \inidx{grammarian}.\footnote{See note~\ref{MLATHRS} on p.~\pageref{MLATHRS} and note \ref{PTJME} on p.~\pageref{PTJME}.} Thus the stanza gives a sense of incongruence with the rest of the work.\footnote{In this regard, all the stanzas at the end of the YVi manuscripts that pay homage to \Patanjali\ (2, 6, 7 and 8) are subject to the same doubt about their authenticity.} Moreover, as mentioned above (pp.~\pageref{patanjaliverseone}\,ff.), this stanza is, if it was a stanza at all, badly transmitted, showing contrast to the immediately preceding and following stanzas that were transmitted with few problems. This gives an impression of different origins. %I am slightly inclined toward the view that this stanza is not genuine part of the text.

 %This further implies that stanzas numbered 1, 2, and 3 all form separate layers.% Such a contrast was one of the things that became inaccessible in the 1952 edition because the orders of stanzas was shuffled.

%determine whether n other words, there is no decisive evidence whether any of the concluding stanzas were by the author himself or someone else.
%It to some degree depends on one's view on the authorship problem. If one considers that the YVi is not a work of the famous \Sankara, then the remaining elements, numbered 1--3 could be part of the original composition of the YVi.\footnote{I exclude number 4, the colophon, in this discussion for now. For, one could argue that the \colophon\ is always independent of the main body of the text. At least, matters are more complicated if we take the placement of the \colophon\ into account. See below.} On the other hand, if one's position is that the YVi is a work of the famous \Sankara, not everything up to 3 (dedicated clearly to the `first' \Sankara) is by himself. The point is that this part, particularly the stanza dedicated to \Sankara, cannot be evidence for or against \Sankara's authorship of the YVi. Rather, the authorship is a factor in determining whether there were more layers in the ending of the YVi.
%Yet I would like to point out that there are some doubts whether the second and third concluding stanzas are by the author of the YVi himself. 

The third stanza\footnote{See p.~\pageref{THETHIRDST}.} is in a similar situation: no clear boundaries around it, but some incongruence with the body of the work. It pays homage to \Sankara\ but in no place in the body of the YVi is he or his \BSBh\ mentioned. Additionally, if we assume that the text of the YVi continued up to this stanza, then the immediately following
%Another possibility is to assume a border just before the colophon and all the concluding stanzas are part of the work, i.e., written by the author of the YVi. In this case, 
colophon that ascribes the YVi to \Sankara\ will contradict what the stanza implies: the author is not \Sankara. The speaker of the stanza cannot be \Sankara\ because he praises \Sankara\ and is from generations later. Then the colophon ascribes the YVi, including this stanza, to \Sankara\ldots ?
%I also have some doubt about the next stanza that salutes \Sankara\ (item numbered 3) belonging to the same layer as the body of the text. 
%I already suspect a gap between the first concluding stanza and the second. Anything that follows is subject to the same degree of suspicion. %If the gap is real, then there are possibilities that the third stanza belongs to the same layer as the second stanza or even further outside layer.

Rather, the stanza and the colophon should be read together.
%On the other hand, the stanza and the colophon are in harmony if we don't consider the stanza to be part of the main text. %can be seen congruent with the following colophon. %, there is nothing odd about it.
%One advantage of viewing the third stanza not coming from the author himself is that we avoid conflict between it and the colophon (item numbered 4) that immediately follows. 
The stanza salutes \Sankara\ and the colophon attributes the work to him. The stanza can simply be understood as paying homage to the author of the text. If we do not wish to explain a contradiction, then the stanza should be placed outside the YVi, coming from the idea that \Sankara\ was the author of the YVi.
 %It is not possible to tell if these two small text portions were introduced by the same person or not, but 
%at least it was not \Sankara\ (no one salutes himself and the person who salutes is from generations later\footnote{See pp.~\pageref{THETHIRDST}\,ff.}) or the author of the YVi himself (the stanza/colophon author thinks the author is \Sankara, but he could not have written the stanza or the colophon). This is when we do not assume a border between the third concluding stanza and the colophon. Anyone can write them as long as the person is not \Sankara\ or the author of the YVi. Besides, It is still not a problem even if \Sankara\ was the author. 
%understanding them as coming from the same notion is more realistic. Note that if we consider that 
%
%The most natural way to read this stanza is alongside the colophon. 
%The third stanza and the colophon probably belong to the same layer in terms of the notion behind them. %Thus, again, my inclination is toward the stanza not being authorial.  %This scenario requires us to postulate a reason why this incongruence was brought about. (Did the mere mention of \Sankara in the stanza evoke the notion that the work was by him?) On the other hand, in the alternative scenario we have to assume a border between the text the author wrote and that which was added during the transmission of the text. This seems a little more likely because I already suspect such a border between the first concluding stanza and the second.

%In the 1952 edition, the order of the dedicatory stanzas was changed and some of the signs of the ending consisting of multiple layers were lost. That was unfortunate.

%has a problem with the \colophon\ that immediately follows . If the original author of the YVi wrote the concluding stanzas up to the \Sankara\ salutation and if the \colophon\ was added by someone else to satisfy convention, \footnote{I highly doubt the users of Sanskrit manuscripts, including those who copied them, considered colophons or sub-colophons authorial. They are often changed, omitted, expanded, or abbreviated more or less by scribes. This is observable in many manuscripts of the same text or sometimes even in apographs. There was little fidelity among scribes when it comes to colophons and sub-colophons.} there is an incongruence between the stanza and the colophon. That is, the stanza salutes \Sankara\ (the author of the stanza is not \Sankara) and the \colophon\ ascribes the text (up to the salutation) to \Sankara. This amounts to the \colophon\ ignoring what the preceding stanza indicates. It is more natural if the \Sankara\ salutation and the \colophon\ belong to the same layer and that they stem from the same notion that the first \Sankara\ was the author of the work.


%Another difficulty I have with seeing stanzas 1--3 as by the same author is the relationship between stanza 3 (dedicated to \Sankara) and the colophon. Let us for the moment assume that the stanza was written by the author of the whole YVi. He salutes \Sankara. Then the colophon ascribes the work (including the salutation) to \Sankara. This is not congruous. Something inexplicable must happen for the colophon (attributing the work to \Sankara) to replace the original colophon that recorded the real author. (There must have been another colophon for the current colophon to be there, and it must be another one for a colophon to be meaningful.) Rather, it is far more plausible that the stanza and the colophon forms a set. The stanza pays homage to the first \Sankara\ in the stanza and the colophon attributes the work to the same \Sankara. It is natural that the \Sankara\ salutation was done because someone thought that the work was by \Sankara. Again, whether the work was truly by \Sankara\ or not, the presence of \Sankara\ salutation is due to the ascription of the work to him. 

Finally the colophon.\footnote{See p.~\pageref{WORKCOL}.} It is about the text (YVi) and is not part of the text. We cannot treat it as in the same level as the body of the text. It is outside of the YVi. 
%The original author might have had some sort of colophon, but what we have is conventional.
%It is, including the phrase \emph{samāptam} \ldots, is just conventional. 
%The format, \emph{govindabhagavatpūjyapādaśiṣyasya} \ldots, is a standard one that ascribes a work to the famous \Sankara. 
%We rarely treat the colophon as by the author of the work. The colophon refers to the text and therefore is by nature from the third party perspective. In fact, as just has been discussed, it appears that the preceding stanza and the colophon come from the same notion that the YVi was by (\emph{that} venerable) \Sankara. This is from a third party perspective. 
%There might have been some sort of colophon in the first ever copy of the YVi, but 
%Regardless of the authorship, it is hard to imagine that the original author had this standard colophon already. %it already had this standard format. The colophon, in its current format, is the least likely part to have been in the original text. %thus from the third party, regardless if it is true. To be thorough, let us consider that the third stanza was in fact by the author of the YVi himself. In that case, the author of the colophon ignored or somehow mistook what the stanza implied and attributed the work to \Sankara. In that case, too, the colophon is not part of the YVi.\footnote{There is another possibility: the colophon, as well as the preceding stanza, was part of the original composition. In this case, the author deliberately attributed his commentary to \Sankara\ by saluting him and attributing the work to \Sankara. This is extremely unlikely. He should have realized that he was saluting himself.} %It is not inconceivable that someone wrote the whole YVi and deliberately attributed it to \Sankara\ by means of the colophon, . I do not think the colophon (item 4), including the phrase ``\emph{samāptam\ldots}'' belongs to the same layer as the main body of the text. The presumed speaker of those statements in the colophon, ``This is the work by \ldots. Here ends the [manuscript of] \ldots,'' should be the scribe. Furthermore, if the author is \Sankara, it is hardly conceivable that he calls himself \emph{bhagavat}. If the author is not \Sankara, then, the colophon is telling un
 
In the end, the first stanza (numbered 1 above) following the conclusion of the commentary is the only stanza that I feel compelled to attribute to the same author who composed the body of the commentary. For the rest, I have some reservations. The likelihood of the subsequent elements being additions during the transmission of the YVi increases as we go down the list. 

Here I would like to call attention to a lack of the word \emph{iti} that was last seen at the end of the commentary and before the first stanza. The word would  clearly mark the end of the text. One of the reasons why there are uncertainties about the extent of the original text, and why the editors of the 1952 edition struggled to make sense out of the arrangement of stanzas, is the lack of that word. The lack may indicate that there was more tampering in the ending than has been discussed here.\footnote{Although I do not have any statistics, that scribal additions being preserved in subsequent copies to such an extent as seen here may already be a rare case. On the one hand, this indicates a complex text transmission process. On the other hand, it might suggest the possibility that even the body of the text itself consists of several layers. If this is the case, it will be much harder to reveal layers than to do so in the ending. Note that at least for the part edited in this volume, I did not consider the possibility of text being tampered by later hands. See pp.~\pageref{what_text}\,ff.} Someone who was adding his own material might have inadvertently removed the clear marking of the end of the text, the word \emph{iti}.
% while adding his material, 
% be a sign of the convoluted situation in the end of the manuscripts may testify to 
%overall impression I get from examining the ending of the YVi is 
%a long period of text transmission of the YVi.
% that jumbled up the ending of the text. I do have sympathies for the editors of the 1952 who felt it necessary to move the stanzas and colophons around.

%\subsection{}

